% capitulos/07-exemplos.tex
% ============================================================
%  EXEMPLOS DE ELEMENTOS GRÁFICOS, TABELAS E CÓDIGOS
%  Este capítulo é apenas demonstrativo e deve ser removido
%  ou adaptado para o conteúdo real da tese/dissertação.
% ============================================================

Este capítulo demonstra os principais elementos visuais disponíveis
neste template, incluindo tabelas simples e complexas, figuras,
subfiguras, fórmulas matemáticas, gráficos estatísticos, blocos de
código-fonte e notas de rodapé.\footnote{As notas de rodapé são
inseridas com o comando \texttt{\textbackslash footnote\{texto\}} e
aparecem automaticamente na parte inferior da página corrente, numeradas
de forma sequencial ao longo de todo o documento.}

% ============================================================
\section{Tabela Simples (Estilo ABNT com booktabs)}
% ============================================================

A \autoref{tab:turmas} exemplifica uma tabela no estilo ABNT,
usando o pacote \texttt{booktabs} com linhas horizontais calibradas
e cabeçalho em negrito, sem linhas verticais.\footnote{O estilo ABNT
não recomenda o uso de linhas verticais em tabelas. Use \texttt{booktabs}
sempre que possível.}

\begin{table}[htb]
  \centering
  \caption{Turmas e participantes analisados na pesquisa.}
  \label{tab:turmas}
  \begin{tabularx}{\textwidth}{lllrr}
    \toprule
    \textbf{Universidade} & \textbf{Turma} & \textbf{Professor}
      & \textbf{Alunos} & \textbf{Respostas} \\
    \midrule
    U1 & U1A-2023 & A & 63  & 3.695 \\
    \thinrule
    U2 & U2A-2023 & B & 13  &   365 \\
    \thinrule
    U2 & U2B-2023 & B & 42  &   794 \\
    \midrule
    U1 & U1A-2024 & A & 28  & 1.266 \\
    \thinrule
    U1 & U1B-2024 & C & 22  &   879 \\
    \thinrule
    U1 & U1C-2024 & D & 34  & 1.206 \\
    \thinrule
    U1 & U1D-2024 & E & 21  &   682 \\
    \midrule
    \multicolumn{2}{l}{2 Universidades / 7 Turmas}
      & 5 Professores & 223 & 8.887 \\
    \bottomrule
  \end{tabularx}
  \fonte{Elaborado pelo autor (2025).}
\end{table}

% ============================================================
\section{Tabela Multicolunas com Divisores Verticais}
% ============================================================

A \autoref{tab:multicoluna} demonstra o uso de divisores verticais
e cabeçalhos agrupados com \texttt{\textbackslash multicolumn},
adequados para tabelas comparativas entre grupos ou períodos.

\begin{table}[htb]
  \centering
  \caption{Respostas corretas e incorretas por questão e por ano.}
  \label{tab:multicoluna}
  \small
  \begin{tabular}{cc r | rr r | rr r}
    \toprule
    \multirow{2}{*}{\textbf{Questão}}
      & \multirow{2}{*}{\textbf{Lista}}
      & \multirow{2}{*}{\textbf{Total}}
      & \multicolumn{3}{c|}{\textbf{2023}}
      & \multicolumn{3}{c}{\textbf{2024}} \\
    \cmidrule(lr){4-6}\cmidrule(lr){7-9}
      & & & \textbf{Corretas} & \textbf{\%}
           & \textbf{Incorretas}
           & \textbf{Corretas} & \textbf{\%}
           & \textbf{Incorretas} \\
    \midrule
    01 & 01 & 1.083 & 104 & 17,22 & 500 & 103 & 21,25 & 376 \\
    \thinrule
    02 & 01 &   612 &  99 & 32,35 & 207 &  98 & 32,03 & 208 \\
    \thinrule
    03 & 01 &   708 &  94 & 25,90 & 269 &  93 & 26,96 & 252 \\
    \thinrule
    04 & 01 &   534 &  89 & 35,32 & 163 &  91 & 32,27 & 191 \\
    \thinrule
    05 & 01 &   503 &  87 & 29,49 & 208 &  96 & 46,15 & 112 \\
    \midrule
    06 & 02 &   492 &  79 & 23,94 & 251 &  80 & 49,38 &  82 \\
    \thinrule
    07 & 02 &   688 &  65 & 18,11 & 294 &  72 & 21,88 & 257 \\
    \thinrule
    08 & 02 &   503 &  70 & 29,29 & 169 &  72 & 27,27 & 192 \\
    \thinrule
    09 & 02 &   404 &  54 & 23,38 & 177 &  65 & 37,57 & 108 \\
    \thinrule
    10 & 03 &   354 &  73 & 34,60 & 138 &  73 & 51,05 &  70 \\
    \bottomrule
  \end{tabular}
  \fonte{Dados da pesquisa (2025).}
\end{table}

% ============================================================
\section{Tabela Longa com Continuação entre Páginas}
% ============================================================

% longtable foi removido do exemplo para reduzir o tempo de compilação
% no Overleaf gratuito. O código abaixo mostra a estrutura para uso
% local ou em instância self-hosted:
%
% \begin{longtable}{p{3.5cm} p{2cm} p{2cm} p{2cm} p{3cm}}
%   \caption{Título da tabela longa.}\label{tab:longa}\\
%   \toprule
%   \textbf{Col A} & \textbf{Col B} & \textbf{Col C}
%     & \textbf{Col D} & \textbf{Col E} \\
%   \midrule
%   \endfirsthead
%   %
%   \multicolumn{5}{r}{\small\itshape Continuação da página anterior}\\
%   \toprule
%   \textbf{Col A} & \textbf{Col B} & \textbf{Col C}
%     & \textbf{Col D} & \textbf{Col E} \\
%   \midrule
%   \endhead
%   %
%   \midrule
%   \multicolumn{5}{r}{\small\itshape Continua na próxima página}\\
%   \endfoot
%   %
%   \bottomrule
%   \endlastfoot
%   %
%   Dado 1 & a & b & c & d \\
%   Dado 2 & e & f & g & h \\
%   % ... repita quantas linhas forem necessárias
% \end{longtable}

% ============================================================
\section{Figura Simples}
% ============================================================

A \autoref{fig:exemplo-simples} ilustra como inserir uma figura
simples centralizada com legenda e fonte abaixo.\footnote{Imagens
devem ser salvas na pasta \texttt{figuras/} e referenciadas pelo
nome sem extensão. Formatos recomendados: PDF (vetorial) ou PNG
(600 dpi mínimo para impressão).}

\begin{figure}[htb]
  \centering
  % Substitua o \fbox abaixo por \includegraphics[width=0.7\textwidth]{nome-arquivo}
  \fbox{\rule{0pt}{5cm}\rule{0.7\textwidth}{0pt}}
  \caption{Perfil de concentração de HgT (ng g$^{-1}$) em função da
           profundidade do solo nos pontos amostrais da ESECAE.
           Barras horizontais indicam o desvio padrão ($n = 30$).}
  \label{fig:exemplo-simples}
  \fonte{Dados da pesquisa (2025).}
\end{figure}

% ============================================================
\section{Figura em Duas Colunas (Subfiguras)}
% ============================================================

A \autoref{fig:subfiguras} demonstra o uso do ambiente
\texttt{subfigure} para apresentar duas figuras lado a lado com
legendas individuais e uma legenda geral.

\begin{figure}[htb]
  \centering
  \begin{subfigure}[b]{0.48\textwidth}
    \centering
    \fbox{\rule{0pt}{4.5cm}\rule{\linewidth}{0pt}}
    \caption{Período seco (julho/2024).}
    \label{fig:subfig-seca}
  \end{subfigure}
  \hfill
  \begin{subfigure}[b]{0.48\textwidth}
    \centering
    \fbox{\rule{0pt}{4.5cm}\rule{\linewidth}{0pt}}
    \caption{Período chuvoso (janeiro/2025).}
    \label{fig:subfig-chuva}
  \end{subfigure}
  \caption{Mapas de concentração de HgT na camada 0--10 cm
           para os dois períodos extremos de amostragem:
           \subref{fig:subfig-seca}~período seco e
           \subref{fig:subfig-chuva}~período chuvoso.}
  \label{fig:subfiguras}
  \fonte{Dados da pesquisa (2025).}
\end{figure}

% ============================================================
\section{Gráfico Boxplot (pgfplots)}
% ============================================================

A \autoref{fig:boxplot} apresenta boxplots comparativos das
concentrações de HgT nos três horizontes de solo e nas três
campanhas de coleta, gerados com o pacote \texttt{pgfplots}.

% ============================================================
\section{Gráfico Boxplot (via R ou Python)}
% ============================================================

O \texttt{pgfplots} não é carregado neste template por ser pesado
demais para o limite de compilação do Overleaf gratuito (20 s).
A abordagem recomendada — e mais prática — é gerar o gráfico
externamente no R ou Python, exportar como \textbf{PDF vetorial}
e inserir com \verb|\includegraphics|, conforme a
\autoref{fig:boxplot}.

\begin{figure}[htb]
  \centering
  % Substitua o \fbox por:
  % \includegraphics[width=\textwidth]{figuras/boxplot_hgt}
  \fbox{\rule{0pt}{6cm}\rule{\textwidth}{0pt}}
  \caption{Distribuição das concentrações de HgT (ng g$^{-1}$)
           por horizonte de solo e campanha de coleta.
           As caixas representam o IQR, a linha central a mediana
           e as hastes os valores extremos não discrepantes.}
  \label{fig:boxplot}
  \fonte{Dados da pesquisa (2025).}
\end{figure}

O \autoref{cod:r-boxplot} mostra o script R para gerar e exportar
o boxplot como PDF pronto para inserção no LaTeX.

\begin{lstlisting}[language=R,
  caption={Script R para gerar boxplot e exportar como PDF.},
  label=cod:r-boxplot]
library(ggplot2)
library(dplyr)

dados <- read.csv("hgt_amostras.csv") |>
  mutate(horizonte = factor(horizonte,
                            levels = c("0-10","10-20","20-30")),
         campanha  = factor(campanha,
                            levels = c("seca","transicao","chuvosa"),
                            labels = c("Seca","Transicao","Chuvosa")))

p <- ggplot(dados, aes(x = horizonte, y = hgt, fill = campanha)) +
  geom_boxplot(alpha = 0.7, outlier.shape = 21, width = 0.6) +
  scale_fill_manual(values = c("#004682","#00783E","#E87722")) +
  labs(x = "Horizonte (cm)",
       y = expression(HgT~(ng~g^{-1}~peso~seco)),
       fill = "Campanha") +
  theme_bw(base_size = 11) +
  theme(legend.position = "top")

# Exporta como PDF vetorial (ideal para LaTeX)
ggsave("figuras/boxplot_hgt.pdf", p,
       width = 16, height = 10, units = "cm")
\end{lstlisting}

% ============================================================
\section{Fórmula Matemática Complexa}
% ============================================================

O modelo de distribuição espacial de HgT foi ajustado pelo
estimador de máxima verossimilhança (MLE) com função de
verossimilhança log-normal multivariada, conforme a
\autoref{eq:mle-completa}:

\begin{equation}
  \label{eq:mle-completa}
  \hat{\boldsymbol{\theta}} =
    \underset{\boldsymbol{\theta}}{\arg\max}
    \left[
      -\frac{n}{2}\ln(2\pi)
      -\frac{1}{2}\ln\left|\boldsymbol{\Sigma}(\boldsymbol{\theta})\right|
      -\frac{1}{2}
        \left(\mathbf{y} - \boldsymbol{\mu}\right)^{\!\top}
        \boldsymbol{\Sigma}^{-1}(\boldsymbol{\theta})
        \left(\mathbf{y} - \boldsymbol{\mu}\right)
    \right]
\end{equation}

\noindent em que $\mathbf{y} \in \mathbb{R}^n$ é o vetor de
concentrações observadas, $\boldsymbol{\mu}$ o vetor de médias,
$\boldsymbol{\Sigma}(\boldsymbol{\theta})$ a matriz de
covariância parametrizada por $\boldsymbol{\theta}$, e
$|\cdot|$ denota o determinante matricial.

A variograma experimental $\hat{\gamma}(h)$ foi estimada por:

\begin{equation}
  \label{eq:variograma}
  \hat{\gamma}(h) =
    \frac{1}{2\,|N(h)|}
    \sum_{(i,j)\,\in\,N(h)}
    \bigl[Z(\mathbf{s}_i) - Z(\mathbf{s}_j)\bigr]^2
\end{equation}

\noindent em que $N(h) = \{(i,j) : \|\mathbf{s}_i - \mathbf{s}_j\| = h\}$
é o conjunto de pares de pontos separados pela distância de lag~$h$,
$|N(h)|$ é a cardinalidade desse conjunto, e $Z(\mathbf{s})$ é
o valor da variável regionalizada na posição $\mathbf{s}$.

O modelo exponencial ajustado à variograma tem a forma:

\begin{equation}
  \label{eq:modelo-exponencial}
  \gamma(h;\, C_0, C_1, a) =
  \begin{cases}
    0 & \text{se } h = 0 \\[4pt]
    C_0 + C_1
      \left[1 - \exp\!\left(-\dfrac{h}{a}\right)\right]
    & \text{se } h > 0
  \end{cases}
\end{equation}

\noindent em que $C_0$ é o efeito pepita (nugget), $C_1$ é o
patamar parcial (sill), e $a$ é o alcance prático (range).

% ============================================================
\section{Exemplo de Código-Fonte (Listings)}
% ============================================================

O \autoref{cod:python-hg} apresenta o script Python utilizado para
o cálculo automatizado do variograma experimental e ajuste do modelo
exponencial às concentrações de HgT.\footnote{O código completo está
disponível no repositório público do projeto em
\url{https://github.com/usuario/ppgca-hg-cerrado}.}

\begin{lstlisting}[language=Python,
  caption={Script Python para calculo do variograma e ajuste de modelo exponencial.},
  label=cod:python-hg]
import numpy as np
from scipy.optimize import curve_fit
from scipy.spatial.distance import pdist, squareform

def variograma_experimental(coords, valores, n_lags=15):
    """
    Calcula o variograma experimental.
    
    Parametros
    ----------
    coords  : np.ndarray, shape (n, 2) - coordenadas (X, Y)
    valores : np.ndarray, shape (n,)   - valores observados
    n_lags  : int - numero de classes de distancia (lags)
    
    Retorno
    -------
    lags, gamma : arrays com distancias e semivariancias medias
    """
    dist_mat = squareform(pdist(coords))
    max_dist = dist_mat.max() / 2
    lag_edges = np.linspace(0, max_dist, n_lags + 1)

    lags, gamma = [], []
    for k in range(n_lags):
        h_min, h_max = lag_edges[k], lag_edges[k + 1]
        mascara = (dist_mat > h_min) & (dist_mat <= h_max)
        i_idx, j_idx = np.where(np.triu(mascara, k=1))
        if len(i_idx) > 0:
            diff = valores[i_idx] - valores[j_idx]
            lags.append((h_min + h_max) / 2)
            gamma.append(np.mean(diff ** 2) / 2)

    return np.array(lags), np.array(gamma)


def modelo_exponencial(h, C0, C1, a):
    """Modelo exponencial de variograma."""
    return C0 + C1 * (1 - np.exp(-h / a))


# --- Uso ---
coords  = np.loadtxt('coords_utm.csv', delimiter=',')
hgt     = np.loadtxt('hgt_ng_g.csv',  delimiter=',')
lags, gama = variograma_experimental(coords, hgt)

popt, _ = curve_fit(modelo_exponencial, lags, gama,
                    p0=[5, 50, 1000], maxfev=5000)
C0_est, C1_est, a_est = popt
print(f'Nugget C0 = {C0_est:.2f}  |  Sill C1 = {C1_est:.2f}'
      f'  |  Range a = {a_est:.1f} m')
\end{lstlisting}

O \autoref{cod:r-anova} apresenta a análise de variância realizada
no ambiente R para comparação das concentrações entre campanhas.

\begin{lstlisting}[language=R,
  caption={Script R para ANOVA de dois fatores e teste post-hoc de Tukey.},
  label=cod:r-anova]
# Carrega pacotes
library(tidyverse)
library(agricolae)

# Le os dados
dados <- read_csv("hgt_amostras.csv") |>
  mutate(campanha  = factor(campanha,
                            levels = c("seca", "transicao", "chuvosa")),
         horizonte = factor(horizonte,
                            levels = c("0-10", "10-20", "20-30")))

# ANOVA de dois fatores
modelo <- aov(hgt ~ campanha * horizonte, data = dados)
summary(modelo)

# Teste post-hoc de Tukey
tukey <- TukeyHSD(modelo)
print(tukey)

# Grafico de interacao
ggplot(dados, aes(x = horizonte, y = hgt, fill = campanha)) +
  geom_boxplot(alpha = 0.7, outlier.shape = 21) +
  scale_fill_manual(values = c("#00783E", "#004682", "#E87722")) +
  labs(x = "Horizonte (cm)",
       y = expression(HgT~(ng~g^{-1})),
       fill = "Campanha") +
  theme_bw(base_size = 12)

ggsave("boxplot_hgt.pdf", width = 18, height = 12, units = "cm")
\end{lstlisting}
